%% kapa: exam=2 round=29 year=2018 period=1 subject=appraisal_practice subject_ko=감정평가실무 seq=4 number=4 points=10
%% source: Q-Net 감정평가사 기출문제 (cst00309) artlSeq=5203398 file=raw/1008/5203398/1교시 감정평가실무.pdf
%% conversion: kordoc 4.24.51 -> tools/gen2.py (auto); review=proofread-sonnet; flags=none
\begin{longquestion}{【문제 4】}
투자자 甲은 1동의 건물 전체를 주거용으로 임대 중인 ‘단지형 연립주택’ 대하여 아래의 < 자료 1 >과 < 자료 2 >를 참고하여, < 자료 3 >과 같이 대상부동산의 가치를 산정하였다. 관련 법규 및 감정평가이론을 참작하고, 제시된 자료를 활용하여 다음의 물음에 답하시오. (10점) \par
\medskip\noindent 물음 1) < 자료 3 >과 같은 甲의 가치산정과정과 결과에 대하여 ‘조소득승수(gross income multiplier)’를 활용하여 점검하되, 최종 조소득승수는 매매사례 (a)와 (b)를 평균하여 산정하시오. (5점) \par
\medskip\noindent 물음 2) 甲의 가치산정 논리에 대하여, 물음 1)의 결과에 기초한 평가 검토의견을 작성하시오. (5점) \par
< 자료 1 > 대상부동산의 개요 \par
1. 총 20개호인 각호의 구조는 2개의 룸으로 구획되어 있고 모두 동일함 \par
2. 보증금 없이 각호 당 월 50만원에 전체를 공실없이 임대중임 \par
< 자료 2 > 인근의 부동산 매매사례 등 \par
1. 매매사례 (a): 총 20개호인 각호의 구조는 3개의 룸으로 구획되어 있고 모두 동일하며, 보증금 없이 각호 당 월 70만원에 전체를 공실없이 임대중이며, 최근 12억원에 거래되었음 \par
2. 매매사례 (b): 총 20개호인 각호의 구조는 4개의 룸으로 구획되어 있고 모두 동일하며, 보증금 없이 각호 당 월 90만원에 전체를 공실없이 임대중이며, 최근 16억원에 거래되었음 \par
< 자료 3 > 甲의 가치산정과정 및 결과 \par
1. 매매사례 (a): 12억원 전체 룸의 수 60개 \par
2. 매매사례 (b): 16억원 전체 룸의 수 80개 \par
3. 룸 당 평균단가: (12억원 + 16억원) / (60룸 + 80룸) = 2,000만원 \par
4. 대상부동산의 시산가치: 2,000만원 × (2룸 × 20개호) = 8억원 \par
\end{longquestion}
